\documentclass{beamer}
\usetheme{Madrid}
\usecolortheme{default}
\usepackage[style=alphabetic]{biblatex}
\addbibresource{ref.bib}
\title[Noisy circuit simulation]
{Pauli path method}
\subtitle{Brief representation}
\author[Siyuan Chen] % (optional, for multiple authors)
{Siyuan Chen\inst{1}}
% \institute[Theoretical Group] % (optional)
% {
 
%   \inst{1}%
%   USTC Shanghai Institute for Advanced Study, Shanghai, China
% }
\date[Representation] % (optional)
{\today}

% \logo{\includegraphics[height=0.8cm]{logo_uoft}}

\definecolor{uoftblue}{RGB}{6,41,88}
\setbeamercolor{titlelike}{bg=uoftblue}
\setbeamerfont{title}{series=\bfseries}
\usepackage{physics}
%%%%%%%%%%%%%%%%%%%%%%%%%%%%%%%%%%%%%%%%%%%%

\begin{document}
\include{preamble.tex}
\frame{\titlepage}

\begin{frame}
\frametitle{Table of Contents}
\tableofcontents
\end{frame}




% \begin{frame}[allowframebreaks]{Pauli path simulation}
%    Consider the depth-$d$ circuit with local depolarizing noise. We want to simulate the average of Pauli operator $O$ with input state $\rho_0$. The state after the circuit writes:
% \begin{align}
%     C(\rho_0) = \mathcal{D}_d\{U_d\cdots \mathcal{D}_1\{ U_1\mathcal{D}_0\{\rho_0\}U_1^\dagger\}\cdots U_d^\dagger \},
% \end{align}
% where $\mathcal{D}$ is the composition of local depolarization noise channel writes $\mathcal{D}^{(j)}(\rho) = (1 - \gamma) \rho + \gamma \tr_j \rho/2 $ . Under Heisenberg picture, the Pauli operator evolve according to $C^{\dagger}(O)$. 
% For local depolarization noise, use the fact that:
% \begin{align}
%     \tr[\mathcal{D}^{(j)}(\rho) P] = \tr[\rho \mathcal{D}^{(j)}(P)],
% \end{align}
% and
% \begin{align}
%     \mathcal{D}^{(j)}(P) = \left\{
%     \begin{aligned}
%        P  &\text{~if~} P^{(j)} = \mathbb{I} \\
%  (1 - \gamma)P  &\text{~if~} P^{(j)} = X,Y,Z \\
%     \end{aligned}
%     \right.
% \end{align}
% Thus the average of operator $O$  is decided by summation:
% \begin{align}
%     \tr[C^\dagger(O)\rho_0] &= \sum_{\vec P = (P_0,\cdots,P_d)} (1 - \gamma)^{\text{wt}(\vec P)} f(\vec P) ,\\
%     \text{where~} f(\vec P) &=\tr(O P_d) \tr(P_d U_d P_{d - 1} U_d^\dagger )\cdots \tr(P_1 U_1 P_{0} U_1^\dagger )  \tr(P_0 \rho_0)
% \end{align}

% \end{frame}




% \begin{frame}[allowframebreaks]{Application: Noisy 2D Clifford + T circuits}
% \begin{definition}[continuous sparse T gate]
% Use $(x,y,t) \in \mathbb{Z}_+^3$ to represent gates implemented on qubit $(x,y)$ at layer $t$.  
% \begin{itemize}
%     \item A subset $A \subset \mathbb{Z}_+^3$ is continuous iff  $\forall x \in A, \exists x \neq y \in A, s.t. y \in \text{Adj}(x) $.
%     \item $T(A)$ denotes the number of T gate in set $A$.
%     \item  A 2D Clifford+T circuit $C$ is $(Q,k)$-sparse if  $\forall A \in \{ A\subset C| |A| \geq k, A \text{~is continous}\}$, $T(A) < Q|A|$. 
% \end{itemize}
% \end{definition}

% \newpage
% \begin{lemma}[Path counting for Continuous Sparse T+Clifford circuit]
%    Given a 2-dimensional, $(Q,k)$-sparse, localized Clifford+T circuit $C$.  
%     \begin{itemize}
%         \item Supposed $\vec P$ has non-vanishing $f_C(\vec P)$,  its support must be continuous.
%         \item For path $\vec P $ with $k \leq \text{wt}(\vec P) \leq w$ ,then $T(\text{supp}[\vec P]) \leq Q w \leq Ql$. 
%         \item Use $N_w$ denotes the number of Pauli paths with non-vanished $f_C(\vec P)$ and $\text{wt}(\vec P) = w$, then $\sum_{w = k}^l N_w \leq 2^{Ql}$.
%     \end{itemize}
% \end{lemma}

% \newpage
% \begin{theorem}
%     Given a $(Q,k)$-sparse 2-dimensional T+Clifford circuit with depolarization noise $\gamma$  and initial product state $\rho_0$, there exist efficient classical algorithm that estimate  $\langle O \rangle$ to precision $\varepsilon $ if $k \sim a \ln n$ and $\gamma > 1 - 1/2^Q$.
% \end{theorem}
% \begin{proof}
%     Take $l > k$, then
%     \begin{align}
%         |\sum_{\text{wt}(\vec P) > l } (1 - \gamma)^{\text{wt}(\vec P)} f(\vec P)| &\leq \sum_{w > l}  N_w (1 - \gamma)^{w}\\
%         &\leq \sum_{w = l}^{nd}  2^{Qw} (1 - \gamma)^{w} \leq \varepsilon
%     \end{align}
% The time complexity related to $\sum_{w =1}^k N_w + \sum_{w=k}^l N_w$, given that each one Pauli path can be calculated in time $O(nd)$.
% \begin{align}
%     \sum_{w =1}^k N_w + \sum_{w=k}^l N_w \leq 2^k + 2^{Ql}
% \end{align}
% \end{proof}
% \begin{align}
%     \mathcal{N}(\rho) &:= (1 - p)\rho + p[\eta_x X\rho X +\eta_y Y\rho Y +\eta_z Z\rho Z].\\
%     X &\to (1- p)X + p(\eta_x - \eta_y - \eta_z),\\
%     I& \to I
% \end{align}
% \end{frame}




\begin{frame}{T-gate + depolarization balance has been done}
    \begin{itemize}
        \item $\langle O \rangle =\sum_w \left[\sum_{|s| = w}f(s)\right] (1 - p)^w.$
        \item As long as $\left[\sum_{|s| = w}f(s)\right] \leq 2^{Qw}$,with $2^Q(1 - p) < 1$, truncation will still works.
        \item   In reference~\cite{gonzalez-garcia_pauli_2024}, they proposed when $T$ gate is "sparse" enough, this condition will be satisfied.
        \item How to related sparseness of T gates with magic resource? There are mainly two choices:
\begin{enumerate}
    \item Gadget-based simulation.
    \item sum-over-Cliffords method.
\end{enumerate}
    \end{itemize}
    
\end{frame}
\begin{frame}{Gadget-based}
\begin{figure}
        \centering
        \includegraphics[width=0.5\linewidth]{pic/gadget_based.png}
        \label{fig:gadget_base}
    \end{figure}
    \begin{itemize}
 \item Only involve state magic, has been well researched.
 \item Move noise forwards to the magic resource state. It can be done efficiently as long as noise channel is mixed by Cliffords.
 \item For depolarization noise, $Z^{a}X^{b} \otimes I \to Z^{a}X^{b} \otimes X^{b}$. How will this channel reduced the magic of state $\ket{0}^{\otimes n} \ket{T}^{\otimes t}$? Stabilizer rank/ stabilizer robustness/Wigner negativity....
 \item New observation: Concentration of gates in third Clifford hierarchy will simplify the simulation of Pauli average.
    \end{itemize}
\end{frame}

\begin{frame}{Gate method}
\begin{itemize}
    \item Is "Gates branch Pauli" a nice protocol to quantify the gate magic and simulation complexity?  
\end{itemize}
    
\end{frame}

\begin{frame}[allowframebreaks]{Reference}
\printbibliography
\end{frame}
\end{document}
